\paragraph{Measurement and response processes} Latent mean comparison requires comparable measurement structure, not merely identical wording \citep{meredith1993,vandenberg2000}. Partial invariance can be useful when a defensible invariant subset anchors comparisons \citep{byrne1989}; freeing problematic parameters does not by itself validate a model. International-survey fit cutoffs depend on group structure and model conditions \citep{rutkowski2014,svetina2017}. Ordinal-invariance and differential-item-functioning approaches offer useful extensions beyond our continuous-moment approximation \citep{millsap2004,svetina2014}. We consequently distinguish directly observed response proportions from claims about an invariant psychological construct.

\paragraph{Overlap, selection and model interpretation} Propensity methods require balance and adequate support as well as assumptions about unmeasured selection \citep{stuart2010,austin2009,austin2015}. Severe distributional separation can make population reweighting unstable; restricting support or using overlap weights defines a different target population \citep{crump2009,li2018}. We do not relabel winsorised DFL ratios as an identified overlap-population effect. SHAP attributions depend on feature dependence and how the background distribution is specified \citep{aas2021,janzing2020}; within-model attribution therefore cannot be promoted to a standardised between-cycle inequality measure.

\paragraph{Digital use, pandemic context and Vietnamese schooling} Reviews distinguish access and duration from use patterns, individual vulnerability and social context \citep{dienlin2020,odgers2020review,odgers2020screen}. Pandemic-era syntheses document distress and isolation concerns but do not identify the source of our PISA cohort differences \citep{racine2021,loades2020,samji2022}. A mixed-method study of Hanoi high-school students reports academic, social and family concerns surrounding anxiety \citep{vietnam2024anxiety}; its local sample and GAD-7 outcome are not interchangeable with national PISA mathematics anxiety. The SMART Schools observational study adds recent evidence on restrictive policies, phone use and well-being \citep{goodyear2025}; policies on paper, enforcement and students' overall use are distinct exposures. These sources motivate alternatives to a uniquely digital ecological explanation, without assigning those alternatives causal effects in our data.

